\section{Conclusions}
The territorial crisis in Santa Cruz de Tenerife is not driven by runaway demographic expansion, but by an artificial regulatory bottleneck:
an obsolete and judicially paralyzed planning framework,
incapable of accommodating even the modest formation of new households, economic modernization, and tertiary transitions.

The current state of urban planning in Santa Cruz de Tenerife is legally fragmented,
given that multiple court rulings have systematically annulled the most recent general plans, as well as their provisional adaptations.

This instability in urban planning generates severe legal uncertainty, particularly regarding building permits and development licenses.
Without an active and up-to-date comprehensive master plan, the legal basis for granting licenses becomes ambiguous.
While some sources point to the partial applicability of the 2005 regulatory framework or transitional provisions,
the judicial suspension of the 2022 municipal ordinance undermined, in particular, the most recent attempt to establish clear land use standards. This situation demands utmost caution regarding any urban development project, as the validity of permits could be challenged due to conflicting legal statuses among planning instruments.
